% GENERATED FILE. Edit canonical lessons, then run npm run generate:books.
% canonical-source-hash: fnv1a64:b817eb0417958567
% canonical-lessons: PT-C495-educar, PT-C495-comportar, PT-C495-castigar, PT-C495-mesada, PT-C495-mimado

\chapter{Educar os filhos}
\label{ch:pt-educar-os-filhos}
\hlchaptermodality{495}

\begin{chapteropening}
\textbf{By the end of this chapter:} \emph{I can talk about bringing up children.}

Complete the last lesson of chapter 495: I can talk about bringing up children.
\end{chapteropening}

\section[educar]{educar — to bring up, to educate}

\label{lesson:PT-C495-educar}



\begin{quote}
Before the new one: say the Portuguese for to crawl, then the Portuguese for to rock.
\end{quote}

\subsection*{You'll want to know: educar}
\textbf{educar} — \textquotedblleft{}to bring up, to educate\textquotedblright{}.

\begin{quote}
Os meus pais educaram-nos com muito carinho.
\end{quote}

\begin{quote}
\emph{My parents brought us up with a lot of love.}
\end{quote}

\emph{Bem-educado} — polite, well brought up.

\begin{cousinweb}
\emph{educar} comes from Latin \emph{ēducāre} \textquotedblleft{}to bring up, to rear\textquotedblright{} (EtymDB).
\end{cousinweb}

\begin{grammarlens}[title={What you've built}]
One more word for telling, explaining and describing at B1.
\end{grammarlens}

\subsection*{Your turn}
\begin{itemize}
\raggedright
  \item \emph{Say it:} \emph{educar}
  \item \emph{Say it:} \emph{educar}, once more
  \item \emph{Say it:} a sentence of your own with \emph{educar}
  \item \emph{Recall:} say the Portuguese for to crawl, then the Portuguese for to rock, then say \emph{educar} again
\end{itemize}

\begin{tcolorbox}[breakable,colback=teal!4,colframe=teal!35!black,title={Before you move on}]
What does educar mean? (\textquotedblleft{}To bring up, to educate\textquotedblright{}.) Say it once more.
\end{tcolorbox}
\section[comportar]{comportar — to behave (comportar-se); to hold}

\label{lesson:PT-C495-comportar}



\begin{quote}
Before the new one: say the Portuguese for to rock, then the Portuguese for to bring up.
\end{quote}

\subsection*{You'll want to know: comportar}
\textbf{comportar} — \textquotedblleft{}to behave (comportar-se); to hold\textquotedblright{}.

\begin{quote}
As crianças comportaram-se bem no restaurante.
\end{quote}

\begin{quote}
\emph{The children behaved well in the restaurant.}
\end{quote}

Reflexive in practice: \emph{comportar-se bem} or \emph{mal}. The noun is \emph{o comportamento}.

\begin{cousinweb}
\emph{comportar} comes from Latin \emph{comportāre} \textquotedblleft{}to carry together\textquotedblright{} (EtymDB), built on \emph{portāre} \textquotedblleft{}to carry\textquotedblright{}.
\end{cousinweb}

\begin{grammarlens}[title={What you've built}]
One more word for telling, explaining and describing at B1.
\end{grammarlens}

\subsection*{Your turn}
\begin{itemize}
\raggedright
  \item \emph{Say it:} \emph{comportar}
  \item \emph{Say it:} \emph{comportar}, once more
  \item \emph{Say it:} a sentence of your own with \emph{comportar}
  \item \emph{Recall:} say the Portuguese for to rock, then the Portuguese for to bring up, then say \emph{comportar} again
\end{itemize}

\begin{tcolorbox}[breakable,colback=teal!4,colframe=teal!35!black,title={Before you move on}]
What does comportar mean? (\textquotedblleft{}To behave (comportar-se); to hold\textquotedblright{}.) Say it once more.
\end{tcolorbox}
\section[castigar]{castigar — to punish}

\label{lesson:PT-C495-castigar}



\begin{quote}
Before the new one: say the Portuguese for to bring up, then the Portuguese for to behave.
\end{quote}

\subsection*{You'll want to know: castigar}
\textbf{castigar} — \textquotedblleft{}to punish\textquotedblright{}.

\begin{quote}
Os pais castigaram o filho: uma semana sem televisão.
\end{quote}

\begin{quote}
\emph{The parents punished their son: a week without television.}
\end{quote}

\emph{Ficar de castigo} — to be grounded.

\begin{cousinweb}
\emph{castigar} is inherited, through Old Portuguese, from Latin \emph{castīgāre} \textquotedblleft{}to correct, to punish\textquotedblright{} (EtymDB).
\end{cousinweb}

\begin{grammarlens}[title={What you've built}]
One more word for telling, explaining and describing at B1.
\end{grammarlens}

\subsection*{Your turn}
\begin{itemize}
\raggedright
  \item \emph{Say it:} \emph{castigar}
  \item \emph{Say it:} \emph{castigar}, once more
  \item \emph{Say it:} a sentence of your own with \emph{castigar}
  \item \emph{Recall:} say the Portuguese for to bring up, then the Portuguese for to behave, then say \emph{castigar} again
\end{itemize}

\begin{tcolorbox}[breakable,colback=teal!4,colframe=teal!35!black,title={Before you move on}]
What does castigar mean? (\textquotedblleft{}To punish\textquotedblright{}.) Say it once more.
\end{tcolorbox}
\section[a mesada]{a mesada — (monthly) pocket money, allowance}

\label{lesson:PT-C495-mesada}



\begin{quote}
Before the new one: say the Portuguese for to behave, then the Portuguese for to punish.
\end{quote}

\subsection*{You'll want to know: a mesada}
\textbf{a mesada} — \textquotedblleft{}(monthly) pocket money, allowance\textquotedblright{}.

\begin{quote}
Recebo uma mesada de vinte euros.
\end{quote}

\begin{quote}
\emph{I get twenty euros pocket money a month.}
\end{quote}

Built on \emph{o mês}: money given each month. Weekly money is \emph{a semanada}.

\begin{cousinweb}
\emph{mesada} is built on \emph{mês} \textquotedblleft{}month\textquotedblright{}, which is inherited from Latin \emph{mēnsis} \textquotedblleft{}month\textquotedblright{} (Lewis \& Short).
\end{cousinweb}

\begin{grammarlens}[title={What you've built}]
One more word for telling, explaining and describing at B1.
\end{grammarlens}

\subsection*{Your turn}
\begin{itemize}
\raggedright
  \item \emph{Say it:} \emph{a mesada}
  \item \emph{Say it:} \emph{a mesada}, once more
  \item \emph{Say it:} a sentence of your own with \emph{a mesada}
  \item \emph{Recall:} say the Portuguese for to behave, then the Portuguese for to punish, then say \emph{a mesada} again
\end{itemize}

\begin{tcolorbox}[breakable,colback=teal!4,colframe=teal!35!black,title={Before you move on}]
What does a mesada mean? (\textquotedblleft{}(monthly) pocket money, allowance\textquotedblright{}.) Say it once more.
\end{tcolorbox}
\section[mimado]{mimado — spoiled}

\label{lesson:PT-C495-mimado}



\begin{quote}
Before the new one: say the Portuguese for to punish, then the Portuguese for (monthly) pocket money.
\end{quote}

\subsection*{You'll want to know: mimado}
\textbf{mimado} — \textquotedblleft{}spoiled\textquotedblright{}.

\begin{quote}
É um menino mimado que quer tudo.
\end{quote}

\begin{quote}
\emph{He is a spoiled little boy who wants everything.}
\end{quote}

Of a child who gets too much: \emph{uma criança mimada}.

\begin{cousinweb}
\emph{mimado} is the participle of \emph{mimar} \textquotedblleft{}to pamper\textquotedblright{}, built on \emph{mimo} \textquotedblleft{}a treat, a caress\textquotedblright{}; the origin of \emph{mimo} is not settled.
\end{cousinweb}

\begin{grammarlens}[title={What you've built}]
One more word for telling, explaining and describing at B1.
\end{grammarlens}

\subsection*{Your turn}
\begin{itemize}
\raggedright
  \item \emph{Say it:} \emph{mimado}
  \item \emph{Say it:} \emph{mimado}, once more
  \item \emph{Say it:} a sentence of your own with \emph{mimado}
  \item \emph{Recall:} say the Portuguese for to punish, then the Portuguese for (monthly) pocket money, then say \emph{mimado} again
\end{itemize}

\begin{tcolorbox}[breakable,colback=teal!4,colframe=teal!35!black,title={Before you move on}]
What does mimado mean? (\textquotedblleft{}Spoiled\textquotedblright{}.) Say it once more.
\end{tcolorbox}
